% Generated from project outputs. Do not edit by hand.
\begin{table}[!htbp]
\centering
\begin{threeparttable}
\caption{Including in-kind by subsample: 2016-2019}
\label{tab:descriptive-2016-2019-subsamples-y2}
\small
\setlength{\tabcolsep}{4pt}
\renewcommand{\arraystretch}{1.15}
\begin{tabularx}{\linewidth}{@{}Xrrrrr@{}}
\toprule
Sample & N & Mean t & SD t & Mean t+1 & SD t+1 \\
\midrule
General & 276 & 2,687.2 & 1,829.6 & 2,753.5 & 1,786.4 \\
Men & 276 & 2,988.3 & 2,201.6 & 3,033.8 & 2,129.6 \\
Women & 276 & 2,173.3 & 1,434.8 & 2,283.4 & 1,512.2 \\
Urban & 276 & 2,961.9 & 1,862.7 & 3,015.5 & 1,804.7 \\
Rural & 271 & 2,003.5 & 1,419.3 & 2,156.1 & 1,569.4 \\
Indigenous & 271 & 1,972.3 & 1,303.3 & 2,058.9 & 1,340.6 \\
Non-Indigenous & 276 & 2,905.3 & 1,879.7 & 2,986.7 & 1,845.5 \\
Bottom 99 percent & 276 & 2,451.8 & 1,362.1 & 2,520.7 & 1,346.0 \\
Bottom 90 percent & 276 & 2,013.5 & 934.9 & 2,070.3 & 913.9 \\
Bottom 50 percent & 272 & 1,063.5 & 426.5 & 1,124.4 & 439.3 \\
\bottomrule
\end{tabularx}
\begin{tablenotes}[flushleft]
\footnotesize
\item Monthly income in quetzales. Unweighted descriptive statistics of survey-weighted cohort means. N counts cohort-transition rows, not individuals or independent cohorts. SD is dispersion across cohort means, not the standard error of a mean. Both incomes must be observed. Subsamples overlap; bottom-income groups are cumulative.
\end{tablenotes}
\end{threeparttable}
\end{table}
